\subsubsection{a. Design of experiment (Latin Hypercube Sampling)}\label{sec:7.2.2.1}

Before any simulation runs, the combinations of parameter values to be tested must
be chosen. A full factorial design is not feasible: ten levels across the nine
parameters of the cable group would give \(10^{9}\) runs. Latin Hypercube Sampling
\cite{mckay1992latin} divides each parameter range into \(N\) intervals of equal
probability and draws one sample from each, then shuffles the assignments across
parameters independently. The property this provides is that every parameter
visits every interval of its own range exactly once, at any \(N\). Coverage of
each parameter individually is therefore guaranteed by construction rather than by
the number of runs, which is what makes a design of a few thousand runs usable
across nine dimensions.

The parameters are listed in Table~\ref{tab:7-1}. Two independent groups were
defined by the presence or absence of boundary cables, giving a seven-parameter
design without them and a nine-parameter design with two, one along each
principal direction. Each group was sampled in two blocks, of \(N = 800\) and
\(N = 1600\), drawn from independent Latin hypercubes with different seeds, so
that the second block extends the first into one pooled space-filling design
rather than duplicating its points. This gives 2400 planned runs per group.

Stratification holds within each block rather than across the pooled design.
Pooling two Latin hypercubes does not yield a single Latin hypercube, but it
preserves even coverage and leaves the parameters close to uncorrelated: the
largest pairwise correlation in the retained designs is 0.062 without cables,
between pressure and the stiffness ratio, and 0.053 with them, between the course
pre-strain and the wale modulus (Figure~\ref{fig:7-11}).

About 55\% of the planned runs reach the surrogate. Three filters apply in
sequence, and Table~\ref{tab:7-3} gives the count at each stage. The first is
numerical convergence to an inflated equilibrium, which 1988 and 1924 runs reach.
The second is a restriction on model validity: configurations with
\(s_{wale}\), \(s_{course} \geq 0.95\) are retained, leaving 1336 and 1314. The
third removes runs whose section profile is too irregular for a section mean
curvature to be meaningful, and costs almost nothing, one run and six. Of the
2400 planned combinations per group, 1335 and 1308 survive, or 56\% and 55\%.

The validity restriction rests on two separate grounds. Below \(s = 0.95\) the
relaxed fabric is larger than the boundary and carries slack, so the simulation
develops compressive membrane stress, which a membrane cannot sustain; a real
fabric wrinkles instead, and the linear orthotropic model of
Section~\ref{sec:5.3.3} does not represent wrinkling. Independently of that, the
biaxial tests from which the constitutive model was derived do not extend into
compression, so those configurations lie outside the domain over which the
material was measured at all.

Attrition is concentrated almost entirely outside the validity box, and is close
to uniform inside it. Of the planned combinations that satisfy
\(s_{wale}, s_{course} \geq 0.95\), 99\% and 97\% converge; of those that do not,
only 62\% and 58\% do. Convergence failures therefore fall where the fabric
cannot be brought into tension, which is the same region the validity
restriction excludes on physical grounds. Within the retained design no marginal
departs detectably from uniform: a Kolmogorov--Smirnov test against the uniform
distribution on the sampled range gives \(p > 0.3\) for every parameter in both
groups, so the retained runs remain a space-filling design over the restricted
box rather than a biased subset of it (Figure~\ref{fig:7-11}).

The indices of Section~\ref{sec:7.2.2.4} therefore describe the response over the
region in which a tensioned membrane solution exists, and not over the full
nominal ranges of Table~\ref{tab:7-1}. This is a restriction on the scope of the
conclusions rather than a defect in the sampling: the excluded region is one in
which the formwork does not function, so the parameters are ranked over the space
in which a design could be built. Because the restriction is a box on the inputs
and not a filter on the outputs, the Saltelli estimator samples the restricted
box directly and never extrapolates into the excluded region.
